\SourceLeaf{66}
$\beta$) Coefficiens, qui initialem excipit, $C^{(1)}$ vocemus, ex scala relationis antea
tradita inveniri potest. Est nimirum
\[
C^{(1)} .  n = -2\int nn\,\mathrm{d}C + nn,
\]
quo substituto et ad calculos revocato, fit
\[
C^{(1)}n = 2\bigl(1-(1-nn)^{\frac{1}{2}}\bigr),
\]
adeoque
\[
C^{(1)} = 2\,\frac{1-\BareRoot{(1-nn)}}{n}.
\]
\par\noindent\hspace*{\parindent}$\gamma$) Ceteros coefficientes ex scala relationis ante evicta una formula generali
circumscribere licet, quae sic enuntiari potest, fore, si ipsius $\log(1+n\cos\varphi)$ in
seriem evolutae coefficiens $h^{\mathrm{tus}} = C^{(h)}$, hunc
\[
= (-1)^{h+1}\,\frac{2}{h}\left(\frac{1-\sqrt{1-nn}}{n}\right)^{h}.
\]
\par\noindent\hspace*{\parindent}Sumamus, legem valere pro $h$ primis, erit tunc ex scala generali, posito $\nu=0$, (7, $\gamma$)
\[
C^{(h+1)} = \frac{-n(h-1)C^{(h-1)}-2h.\,C^{(h)}}{n(h+1)}, 
\]
fit substituendo debitos valores
\[
C^{(h+1)} = -\frac{\begin{gathered}n.(h-1)(-1)^{h}\dfrac{2}{h-1}\left(\dfrac{1-\sqrt{1-nn}}{n}\right)^{h-1}\\{}+2h(-1)^{h+1}\dfrac{2}{h}\left(\dfrac{1-\BareRoot{1-nn}}{n}\right)^{h}\end{gathered}}{n.(h+1)},
\]
sive, separato factore communi
\[
C^{(h+1)} = -\frac{(-1)^{h}2\left(\dfrac{1-\sqrt{1-nn}}{n}\right)^{h-1}}{n(h+1)}\left(n-\frac{2(1-\sqrt{1-nn})}{n}\right).
\]
Est vero factor uncinis inclusus
\[
= \frac{nn-2+2\sqrt{1-nn}}{n} = -\frac{(1-\BareRoot{(1-nn)})^{2}}{n}.
\]
Si igitur alter per illum revera multiplicetur, fit
\[
C^{(h+1)} = (-1)^{h+2}\,\frac{2}{h+1}\left(\frac{1-\BareRoot{(1-nn)}}{n}\right)^{h+1},
\]
unde perspicitur legis generalitas.

